Die Evaluation der barrierefreien Umsetzung erfolgte kontinuierlich und begleitend während der gesamten Entwicklungsphase. Dabei kamen dieselben Test-Tools zum Einsatz, die bereits in der Analyse bestehender Webanwendungen verwendet wurden - darunter unter anderem axe DevTools und WAVE. Diese automatisierten Tests ermöglichten eine schnelle Rückmeldung über häufige Barrierefreiheitsprobleme wie Kontraste, fehlende ARIA-Attribute oder semantische Fehler und wurden direkt in den Entwicklungsprozess integriert. Diese Tools gaben am Ende keine Fehler mehr aus.

Deswegen fanden zusätzlich manuelle Tests statt, um die Qualität der Barrierefreiheit über rein technische Kriterien hinaus zu überprüfen. Dazu zählten etwa die Navigation per Tastatur, das gezielte Setzen und Sichtbarmachen von Fokuszuständen sowie der Einsatz von Screenreadern. So wurden die einzelnen Komponenten genauer unter die Lupe genommen und beobachtet, wie sie sich mit anderen Interaktionen wie der Maus und dem Auge verhalten.

Das erwies sich als besonders hilfreich z.B. bei dem Account-Dropdown. So konnte mit der Tastatur sichergestellt werden, dass das Untermenü aufrufbar ist und die einzelnen Punkte richtig vom Screenreader als Links ausgegeben werden. Genau so bei den Helfer Texten der Input Felder. Hier haben die Test Tools alle keine Fehler ausgegeben. Aber wenn man sich mit dem Screenreader, durch Fokussieren des Input Feldes, die Helfer Texte ausgeben lassen möchte, wird nichts ausgegeben. Erst durch die Verknüpfung der beiden Elemente, durch ARIA, wurden die Helfer Texte vom Screenreader beim Fokussieren berücksichtigt.

Diese Kombination aus automatisiertem und manuellem Testing erwies sich als effektiv, um frühzeitig Schwachstellen zu identifizieren und gezielt zu optimieren. Insbesondere bei interaktiven Komponenten wie Formularen, Navigationselementen oder dem hinzugefügtem FAQ-Bereich und Account-Dropdown konnte so sichergestellt werden, dass nicht nur technische Kriterien erfüllt sind, sondern auch die tatsächliche Nutzungserfahrung barrierefrei gestaltet ist.
